
\exo

On note $\cD$ la droite d'équation $2x-3y+5=0$, et $\cD'$ la parallèle à $\cD$ passant par le point $A\couple{-2}{2}$.
\begin{enumerate}
\item Justifier que le vecteur $\vecu\couple{3}{2}$ est directeur de $\cD'$.
\item Déterminer une équation cartésienne de $\cD'$.
\end{enumerate}

\exo

On note $\cD$ la droite d'équation $y=-5x+3$, et $\cD'$ la parallèle à $\cD$ passant par le point $A\couple{3}{8}$.
\begin{enumerate}
\item Justifier que le coefficient directeur de $\cD'$ est égal à $-5$.
\item Déterminer une équation réduite de $\cD'$.
\end{enumerate}

\exo

On considère les points $A\couple{3}{2}$, $B\couple{4}{2}$ et $C\couple{3}{5}$.
\begin{enumerate}
\item Déterminer une équation de la droite $\cD$, parallèle à $(AB)$ et passant par $C$.
\item Déterminer une équation de la droite $\cD'$, parallèle à $(BC)$ et passant par $A$.
\item Déterminer une équation de la droite $\cD''$, parallèle à $(AC)$ et passant par $B$.
\end{enumerate}

\exo

Dans chaque cas, déterminer une équation de la droite $\cD'$, parallèle à $\cD$ et passant par $A$.
\vspace{-6pt}
\setlength{\columnsep}{1cm}
\setlength{\columnseprule}{0pt}
\begin{multicols}{2}
\begin{enumerate}
\item $\cD$ : $4x+2y-5=0$ et $A\couple{1}{2}$.
\medskip
\item $\cD$ : $y=-2x+7$ et $A\couple{4}{1}$.
\item $\cD$ : $x+4=0$ et $A\couple{11}{2}$.
\medskip
\item $\cD$ : $x-2y+7=0$ et $A\couple{5}{-2}$.
\end{enumerate}
\end{multicols}
\vspace{-8pt}

\exo

On considère le système $(S)~:~\left\{\setlength\arraycolsep{2pt}\begin{array}{rcrcl}
x & - & 6y & = & 70 \\
6x & - & y & = & 70
\end{array}\right.$

Les affirmations suivantes sont-elles vraies ou fausses ? Justifier la réponse.
\smallskip
\begin{enumerate}
\item Le couple $\couple{10}{-10}$ est une solution de $(S)$.
\item Le couple $\couple{-10}{10}$ n'est pas une solution de $(S)$.
\item Le couple $\couple{10}{-10}$ est la solution de $(S)$.
\end{enumerate}

\exo

Résoudre chaque système par substitution.
\vspace{-6pt}
\setlength{\columnsep}{1cm}
\setlength{\columnseprule}{0pt}
\begin{multicols}{3}
\begin{enumerate}
\item $\left\{\setlength\arraycolsep{2pt}\begin{array}{rcrcl}
3x & + & y & = & -1 \\
5x & - & 4y & = & -13
\end{array}\right.$
\medskip
\item $\left\{\setlength\arraycolsep{2pt}\begin{array}{rcrcl}
-x & + & 5y & = & -7 \\
10x & + & 3y & = & 17
\end{array}\right.$
\item $\left\{\setlength\arraycolsep{2pt}\begin{array}{rcrcrcl}
2x & - & y & + & 1 & = & 0 \\
-3x & + & 4y & + & 1 & = & 0
\end{array}\right.$
\medskip
\item $\left\{\setlength\arraycolsep{2pt}\begin{array}{rcrcrcl}
x & - & 3y & + & 4 & = & 0 \\
2x & - & 5y & + & 1 & = & 0
\end{array}\right.$
\item $\left\{\setlength\arraycolsep{2pt}\begin{array}{rcrcrcl}
-3x & + & 2y & - & 1 & = & 0 \\
5x & + & y & + & 2 & = & 0
\end{array}\right.$
\medskip
\item $\left\{\setlength\arraycolsep{2pt}\begin{array}{rcrcrcl}
-x & + & 5y &  &  & = & 0 \\
3x & - & 4y & + & 1 & = & 0
\end{array}\right.$
\end{enumerate}
\end{multicols}
\vspace{-6pt}

\exo

Résoudre chaque système par combinaison.
\vspace{-6pt}
\setlength{\columnsep}{1cm}
\setlength{\columnseprule}{0pt}
\begin{multicols}{3}
\begin{enumerate}
\item $\left\{\setlength\arraycolsep{2pt}\begin{array}{rcrcl}
8x & - & 5y & = & -7 \\
x & + & 5y & = & 16
\end{array}\right.$
\medskip
\item $\left\{\setlength\arraycolsep{2pt}\begin{array}{rcrcl}
3x & + & 5y & = & 16 \\
3x & + & 3y & = & 12
\end{array}\right.$
\item $\left\{\setlength\arraycolsep{2pt}\begin{array}{rcrcl}
3x & + & 5y & = & 31 \\
5x & + & 4y & = & 43
\end{array}\right.$
\medskip
\item $\left\{\setlength\arraycolsep{2pt}\begin{array}{rcrcl}
7x & + & 4y & = & 1 \\
3x & + & 5y & = & 7
\end{array}\right.$
\item $\left\{\setlength\arraycolsep{2pt}\begin{array}{rcrcrcl}
-2x & + & 5y &  &  & = & 0 \\
3x & - & 4y & - & 3 & = & 0
\end{array}\right.$
\medskip
\item $\left\{\setlength\arraycolsep{2pt}\begin{array}{rcrcrcl}
5x & + & 3y & - & 1 & = & 0 \\
-7x & - & 2y & + & 3 & = & 0
\end{array}\right.$
\end{enumerate}
\end{multicols}
\vspace{-6pt}

\exo

Résoudre chaque système en utilisant la méthode la plus adaptée.
\vspace{-6pt}
\setlength{\columnsep}{1cm}
\setlength{\columnseprule}{0pt}
\begin{multicols}{3}
\begin{enumerate}
\item $\left\{\setlength\arraycolsep{2pt}\begin{array}{rcrcl}
2x & - & 3y & = & 1 \\
3x & + & 5y & = & 11
\end{array}\right.$
\item $\left\{\setlength\arraycolsep{2pt}\begin{array}{rcrcl}
3x & + & y & = & 1 \\
4x & - & 3y & = & 10
\end{array}\right.$
\item $\left\{\setlength\arraycolsep{2pt}\begin{array}{rcrcrcl}
2x & + & 3y & + & 5 & = & 0 \\
3x & + & 2y & - & 25 & = & 0
\end{array}\right.$
\end{enumerate}
\end{multicols}
\vspace{-6pt}

\exo

Déterminer l'intersection des deux droites $\cD$ et $\cD'$ dont on donne les équations.
\vspace{-6pt}
\setlength{\columnsep}{1cm}
\setlength{\columnseprule}{0pt}
\begin{multicols}{2}
\begin{enumerate}
\item $\cD$ : $2x-3y-1=0$ et $\cD'$ : $-4x+3y+2=0$.
\medskip
\item $\cD$ : $-3x+2y+1=0$ et $\cD'$ : $x+3y-3=0$.
\item $\cD$ : $x-y+1=0$ et $\cD'$ : $-3x+3y-2=0$.
\medskip
\item $\cD$ : $2x3y+1=0$ et $\cD'$ : $-6x+3y-3=0$.
\end{enumerate}
\end{multicols}
\vspace{-8pt}

\exo

\vspace{-6pt}
\setlength{\columnsep}{1cm}
\setlength{\columnseprule}{0pt}
\begin{multicols}{2}
\raggedcolumns
Déterminer les coordonnées du point d'intersection des droites $\cD_1$ et $\cD_2$ représentées ci-contre.

\begin{center}
\def\xmin{-2}
\def\xmax{10}
\def\ymin{-3}
\def\ymax{6}
\def\xunite{0.6cm}
\def\yunite{0.6cm}
\psset{unit=\xunite, xunit=\xunite, yunit=\yunite, algebraic=true, dimen=middle, dotstyle=*, dotsize=3pt 0, linewidth=0.8pt, arrowsize=3pt 2, arrowinset=0.25, comma=true}
\begin{pspicture*}(\xmin,\ymin)(\xmax,\ymax)
\psgrid[xunit=\xunite, yunit=\yunite, subgriddiv=0, gridlabels=0, gridcolor=lightgray](0,0)(\xmin,\ymin)(\xmax,\ymax)
\psaxes[showorigin=false,labelFontSize=\scriptstyle, xAxis=true, yAxis=true, Dx=1, Dy=1, ticksize=-2pt 0, subticks=1]{->}(0,0)(-1.99,-2.99)(\xmax,\ymax)[{\footnotesize $x$},135][{\footnotesize $y$},-45]
\psplot[linewidth=1.2pt, plotpoints=200]{\xmin}{\xmax}{2/3*(x-2)-1}
\psplot[linewidth=1.2pt, plotpoints=200]{\xmin}{\xmax}{0.25*x+3}
\begin{footnotesize}
\uput{4pt}[315](8,3){$\cD_1$}
\uput{4pt}[110](8,5){$\cD_2$}
\end{footnotesize}
\end{pspicture*}
\end{center}

\end{multicols}
\vspace{-6pt}

\exo

Dans chaque cas, déterminer les nombres $a$ et $b$ de sorte que la condition soit vérifiée.
\medskip
\begin{enumerate}
\item Le couple $\couple{2}{a}$ est solution du système $\left\{\setlength\arraycolsep{2pt}\begin{array}{rcrcl}
2x & - & 3y & = & 13 \\
5x & + & by & = & -14
\end{array}\right.$.
\item Le couple $\couple{5}{-2}$ est solution du système $\left\{\setlength\arraycolsep{2pt}\begin{array}{rcrcl}
ax & + & 2y & = & 1 \\
3x & + & by & = & -1
\end{array}\right.$.
\end{enumerate}

\exo

Une salle de spectacle propose des pièces de théâtre et des concerts.

Pour $170$ \euro{}, Alice a réservé deux places pour une pièce de théâtre, et quatre pour un concert.

Pour $135$ \euro{}, Bob a réservé trois places pour une pièce de théâtre, et deux pour un concert.

Déterminer le prix d'une place pour une pièce de théâtre, et le prix d'une place pour un concert.

\exo

En prenant $14$ douches et $21$ bains chaque semaine, une famille de cinq personnes consomme \np{5040}~litres d'eau. Avec $28$ douches et $7$ bains, elle consommerait seulement \np{3080}~litres d'eau.

Quelle quantité d'eau consomme-t-on dans cette famille pour une douche ? Et pour un bain ?

\exo

Alice possède $20$ pièces de monnaie dans sa tirelire, de $1$ ou $2$ euros chacune, pour un total de $36$ euros.

Combien de pièces de chaque sorte possède-t-elle ?

\exo

Bob souhaite acheter un téléviseur et une tablette.

Dans un premier magasin, ces deux articles lui coûteraient $714$ \euro{}.

Dans un second magasin, le téléviseur bénéficie d'une remise de \SI[retain-explicit-plus]{30}{\percent} et la tablette est vendue $50$ \euro{} moins cher. Le montant total serait alors de $535$ \euro{}.

On note $x$ et $y$ les prix respectifs du téléviseur et de la tablette dans le premier magasin.
\begin{enumerate}
\item Traduire les données du problème par un système de deux équations d'inconnues $x$ et $y$.
\item Déterminer $x$ et $y$.
\end{enumerate}

\exo

De la ville A à la ville B, Alice roule sur autoroute à une vitesse constante. Si elle augmentait sa vitesse de \SI[per-mode=symbol]{20}{\km\per\hour}, elle gagnerait $6$ minutes mais si elle diminuait sa vitesse de \SI[per-mode=symbol]{20}{\km\per\hour}, elle perdrait $9$ minutes.

Quelle est la distance de la ville A à la ville B ?

